\documentclass[../../../thesis.tex]{subfiles}
\begin{document}
\subsection{Štylistika}
Niektorí oponenti vyčítajú študentom príliš dlhé súvetia,
iní zas príliš krátke.
Pravda je, že jednoduché vety pôsobia školácky,
zatiaľ čo dlhé súvetia sú často nezrozumiteľné a~únavné.

V prvom rade sa snažíme nevrstviť podraďovacie súvetia.
Vo vete \emph{Elektrostatické pole je fyzikálne pole,
ktoré tvoria elektrické náboje, ktoré sú v pokoji}
je dvakrát použitá spojka ktoré,
čo je síce prípustné, avšak nie príliš estetické.
Vetu môžeme opraviť takto:
\emph{Elektrostatické pole je fyzikálne
pole tvorené elektrickými nábojmi, ktoré sú v pokoji.}
Ak sa chceme vyhnúť trpnému rodu,
môžeme vetu preformulovať nasledujúcim spôsobom:
\emph{Elektrostatické pole tvoria elektrické náboje,
ktoré sú v~pokoji.}
Vypadol síce pojem fyzikálne pole,
ale zmysel vety zostal nezmenený.

Správne a~plynulo bude veta vyzerať aj v~tomto tvare:
\emph{Elektrostatické pole je fyzikálne pole,
ktoré tvoria elektrické náboje v pokoji.}
V prípade potreby môžeme vetu napísať aj inak:
\emph{Fyzikálne pole elektrických nábojov,
ktoré sú v pokoji, nazývame elektrostatické pole.}

Obmieňame štruktúru po sebe nasledujúcich viet:
\emph{Z výsledkov merania je zrejmé,
že predpoklad o~zvyšovaní pohyblivosti nosičov náboja
s~teplotou bol správny.
Na začiatku práce sme hovorili o~tom,
že toto tvrdenie podporíme hodnovernými experimentálnymi dátami.}
Obe súvetia sú podraďovacie so spojkou že.
Aby sme sa vyhli opakovaniu rovnakého typu viet,
môžeme prvú vetu prepísať:
\emph{Výsledky merania potvrdili predpoklad o~zvyšovaní
pohyblivosti nosičov náboja s~rastúcou teplotou.}
Druhú vetu ponecháme bez zmeny.

Veľmi osviežujúco pôsobí, ak medzi dlhé a~kvetnaté súvetia
občas vložíme jednoduchú holú vetu.
Použijeme predchádzajúci príklad:
\emph{Na začiatku práce sme hovorili o~tom,
že predpoklad o~zvyšovaní pohyblivosti nosičov náboja
s~rastúcou teplotou podporíme hodnovernými
experimentálnymi dátami.
Merania ho potvrdili.}
Tento malý trik je nečakane účinný a~prispieva k~lepšiemu
toku myšlienok.

Pozor, v~texte pozostávajúcom z~krátkych jednoduchých viet
je niekoľkoriadkové súvetie desivé:
\emph{Pohyblivosť rastie s~teplotou.
Hovorili sme o~tom už na začiatku.
Tvrdenie ešte podporíme experimentom.
Ukazuje sa, že sme predpoklad o~rastúcej pohyblivosti
nosičov náboja so zvyšujúcou sa teplotou,
pokiaľ berieme do úvahy výsledky meraní,
formulovali správne.}

Aby bol písaný text zaujímavý a~udržal čitateľov záujem,
používame stredne dlhé súветia pozostávajúce maximálne
z~dvoch až troch viet.
Občas text oživíme jednoduchou krátkou vetou.
Dávame si pri tom pozor,
aby táto činnosť nebola príliš schematická.
\end{document}